\section*{ЗАКЛЮЧЕНИЕ}
\addcontentsline{toc}{section}{\protect ЗАКЛЮЧЕНИЕ}

В результате работы была изучена архитектура Mamba применительно к задаче нейронного машинного перевода.
Были рассмотрены две постановки sequence-to-sequence задачи: классическая схема encoder-decoder и модель только с декодером, где исходное предложение и перевод образуют одну авторегрессионную последовательность.
Также была кратко описана связь Mamba с моделями пространства состояний и идея избирательного обновления скрытого состояния.

В практической части были обучены и сравнены четыре модели: LSTM с механизмом внимания, Transformer, Mamba encoder-decoder и Mamba decoder-only.
Основной метрикой качества была BLEU, дополнительно учитывалось число параметров.
На русско-английском корпусе OPUS-100 Mamba в схеме encoder-decoder показала качество, близкое к Transformer, и превзошла LSTM.

Модель Mamba только с декодером уступила encoder-decoder вариантам по BLEU.
Это показывает, что для данной задачи явное кодирование исходного предложения остается полезным.
В то же время decoder-only схема работоспособна и может рассматриваться как отдельное направление для дальнейшего улучшения Mamba-моделей машинного перевода.
